% -*- coding: utf-8 -*-
\chapter{AIによるコード生成とテスト}
\label{chap6}

\begin{chapterbox}
\textbf{【この章で学ぶこと】}
\begin{itemize}
\item \textbf{AIを活用したコード生成の基本フローとプロンプト設計を習得する} --- 自然言語からコードを生成するプロセスを理解し、効果的なプロンプトを書けるようになる
\item \textbf{テストコードの自動生成とAIコードレビューを実践できる} --- テスト駆動開発の考え方を学び、AIを活用した品質向上手法を身につける
\item \textbf{AI生成コードの品質・セキュリティリスクを理解し対処できる} --- 生成コードを盲信せず、適切に検証・修正する力を養う
\end{itemize}
\end{chapterbox}
\vspace{5mm}

%======================================================================
\section{AIとソフトウェア開発の現在}
\label{sec:6.1}
\index{こーどせいせいAI@コード生成AI}

\subsection{コード生成AIの普及状況}

2024年以降、AIによるコード生成は研究者・開発者の日常的なツールとなった。主要なツールとその特徴を以下に整理する。

\begin{table}[htbp]
\centering
\caption{主要なコード生成AIツール}
\label{tab:code-ai-tools}
\footnotesize
\begin{tabular}{lp{4cm}p{3.5cm}}
\toprule
ツール & 特徴 & 料金 \\
\midrule
GitHub Copilot & VS Code統合、コード補完に特化 & 学生無料（GitHub Education） \\
Cursor & AI統合型エディタ、対話的開発 & 無料枠あり、Pro月\$20 \\
ChatGPT / GPT-4 & 汎用的、説明付きコード生成 & 無料枠あり、Plus月\$20 \\
Gemini Pro & Google連携、長文コンテキスト & 無料枠あり \\
Claude & 長文コンテキスト、正確な推論 & 無料枠あり、Pro月\$20 \\
\bottomrule
\end{tabular}
\end{table}

\subsection{開発者のAI利用統計}

GitHubの2024年開発者調査によると、回答者の92\%がAIコーディングツールを職場内外で利用しており、70\%以上が「生産性が向上した」と回答している。特に以下の場面での活用が報告されている。

\begin{itemize}
\item \textbf{定型的なコードの作成}（ボイラープレートコード）: 87\%
\item \textbf{新しいプログラミング言語の学習}: 73\%
\item \textbf{バグの特定と修正}: 68\%
\item \textbf{テストコードの作成}: 62\%
\end{itemize}

\subsection{AIが変える開発ワークフロー}

従来の開発ワークフローでは、プログラマが仕様を読み、設計し、コードを一行ずつ書いていた。AIの登場により、このプロセスは大きく変わりつつある。

\textbf{従来のワークフロー:}
\begin{verbatim}
仕様理解 → 設計 → コーディング → テスト → デバッグ → レビュー
\end{verbatim}

\textbf{AI活用ワークフロー:}
\begin{verbatim}
仕様記述（プロンプト） → AI生成 → 人間がレビュー・修正
    → AIでテスト生成 → 検証
\end{verbatim}

重要なのは、AIは「コードを書く」部分を加速するが、「何を作るか」「なぜ作るか」を考えるのは依然として人間の仕事であるということだ。大学院生にとって、研究のロジックを明確にする力がますます重要になっている。

%======================================================================
\section{AIによるコード生成の基本}
\label{sec:6.2}
\index{こーどせいせい@コード生成}

\subsection{基本フロー}

AIによるコード生成は、以下の3ステップで行う。

\begin{verbatim}
Step 1: 自然言語で仕様を記述する（プロンプト作成）
Step 2: AIがコードを生成する
Step 3: 人間がレビュー・テスト・修正する
\end{verbatim}

このうち、最も重要なのは\textbf{Step 1（プロンプト作成）}である。指示が曖昧だと、生成されるコードも意図に沿わないものになりやすい。明確で具体的な仕様を書くことが、良いコードを得る鍵である。

\subsection{効果的なプロンプトの4要素}

AIに高品質なコードを生成させるためには、以下の4つの要素をプロンプトに含めることが重要である。

\begin{table}[htbp]
\centering
\caption{効果的なプロンプトの4要素}
\label{tab:prompt-elements}
\resizebox{\textwidth}{!}{%
\begin{tabular}{lll}
\toprule
要素 & 内容 & 例 \\
\midrule
1. 言語・フレームワーク指定 & 使用する言語やライブラリを明示 & 「Python 3.10、pandasを使用」 \\
2. 入出力仕様 & 何を受け取り、何を返すか & 「CSVファイルパスを受け取り、 \\
 & & DataFrameを返す」 \\
3. 制約条件 & パフォーマンス、エラー処理等 & 「メモリ効率を重視、 \\
 & & 欠損値はNaNで保持」 \\
4. 例示 & 具体的な入出力例 & 「入力: 'data.csv' → \\
 & & 出力: 100行×5列のDataFrame」 \\
\bottomrule
\end{tabular}}
\end{table}

\subsection{Bad prompt vs Good prompt の比較}

\textbf{Bad prompt（悪い例）:}

\begin{promptbox}
CSVを読み込む関数を書いて
\end{promptbox}

この指示では、言語もライブラリも不明であり、エラー処理の有無も分からない。AIは何らかのコードを返すが、研究の要件に合うかは運次第である。

\textbf{Good prompt（良い例）:}

\begin{promptbox}
以下の仕様でPython関数を作成してください。

\#\# 関数名
load\_experiment\_data

\#\# 目的
実験結果のCSVファイルを読み込み、前処理済みのDataFrameを返す

\#\# 使用ライブラリ
- Python 3.10
- pandas 2.0以上

\#\# 入力
- file\_path: str --- CSVファイルのパス

\#\# 出力
- pd.DataFrame --- 以下の前処理を適用済み
  1. 列名を英語の小文字に統一（スペースはアンダースコアに変換）
  2. 数値列の欠損値は列の中央値で補完
  3. 重複行を削除

\#\# エラー処理
- ファイルが存在しない場合: FileNotFoundError を発生
- CSVの解析に失敗した場合: ValueError を発生（メッセージにファイル名を含む）
\end{promptbox}

このように具体的な仕様を書くことで、AIはより正確で実用的なコードを生成できる。

\textbf{AIが生成するコードの例:}

\begin{lstlisting}[language=Python]
import pandas as pd
import os


def load_experiment_data(file_path: str) -> pd.DataFrame:
    """
    実験結果のCSVファイルを読み込み、前処理済みのDataFrameを返す。

    Parameters
    ----------
    file_path : str
        CSVファイルのパス

    Returns
    -------
    pd.DataFrame
        前処理済みのDataFrame

    Raises
    ------
    FileNotFoundError
        ファイルが存在しない場合
    ValueError
        CSVの解析に失敗した場合
    """
    # ファイル存在確認
    if not os.path.exists(file_path):
        raise FileNotFoundError(
            f"ファイルが見つかりません: {file_path}"
        )

    # CSV読み込み
    try:
        df = pd.read_csv(file_path)
    except Exception as e:
        raise ValueError(
            f"CSVの解析に失敗しました ({file_path}): {e}"
        )

    # 1. 列名を英語の小文字に統一
    df.columns = (
        df.columns
        .str.strip()
        .str.lower()
        .str.replace(" ", "_")
    )

    # 2. 数値列の欠損値を中央値で補完
    numeric_cols = df.select_dtypes(include="number").columns
    for col in numeric_cols:
        median_val = df[col].median()
        df[col] = df[col].fillna(median_val)

    # 3. 重複行を削除
    df = df.drop_duplicates().reset_index(drop=True)

    return df
\end{lstlisting}

%======================================================================
\section{主要ツールの活用法}
\label{sec:6.3}

\subsection{Gemini Pro / ChatGPT / Claude でのコード生成}

ブラウザベースのAIチャットツールは、最も手軽にコード生成を始められる方法である。

\textbf{活用のコツ:}

\begin{enumerate}
\item \textbf{新しいチャットで始める} --- 過去の文脈が混ざらないよう、コード生成タスクは新しいセッションで行う
\item \textbf{コンテキストを最初に提供する} --- 「私は大学院生で、心理学実験のデータを分析しています」のように背景を伝える
\item \textbf{段階的に依頼する} --- 一度に全てを求めず、まず骨格を作り、次にエラー処理を追加するなど段階的に進める
\end{enumerate}

\textbf{プロンプト例（Claude）:}

\begin{promptbox}
私は神経科学の大学院生です。脳波（EEG）データの解析スクリプトを
作成しています。

以下の関数を Python で作成してください。
- MNE-Python ライブラリを使用
- .fif 形式のEEGデータを読み込む
- バンドパスフィルタ（1-40Hz）を適用
- アーティファクト除去（閾値: ±100μV）を行う
- エポック分割（-0.2〜0.8秒、イベントID指定可能）を行う

各処理ステップにログ出力を含めてください。
\end{promptbox}

\subsection{GitHub Copilot の活用}
\index{GitHub Copilot@GitHub Copilot}

GitHub Copilotは、VS Code上でリアルタイムにコード補完を行うツールである。\textbf{大学院生はGitHub Educationを通じて無料で利用できる}。

\textbf{セットアップ手順:}

\begin{enumerate}
\item GitHub Education（\texttt{https://education.github.com}）で学生認証を行う
\item VS Codeの拡張機能から「GitHub Copilot」をインストール
\item GitHubアカウントでサインイン
\item \texttt{.py}ファイルを開いてコメントを書くと、自動的にコード候補が表示される
\end{enumerate}

\textbf{効果的な使い方:}

\begin{lstlisting}[language=Python]
# --- コメントで意図を書くと、Copilotが補完する ---

# CSVファイルを読み込み、欠損値の数を列ごとに表示する関数
def show_missing_values(file_path):
    # ← ここでTabキーを押すと、Copilotがコードを提案する
\end{lstlisting}

\textbf{Tips:}
\begin{itemize}
\item コメントは英語でも日本語でも良いが、具体的に書くほど精度が上がる
\item \texttt{Alt + ]} で次の候補、\texttt{Alt + [} で前の候補を表示できる
\item Copilot Chat（サイドパネル）を使えば、対話的にコードを改善できる
\end{itemize}

\subsection{対話的コード開発のテクニック}
\index{たいわてきこーどかいはつ@対話的コード開発}

AIとの対話を通じてコードを段階的に改善する方法を、「対話的コード開発」と呼ぶ。以下にその流れを示す。

\textbf{Step 1: 最初のプロンプト}
\begin{promptbox}
二群の平均値を比較する統計検定を行うPython関数を書いてください。
scipy.stats を使用してください。
\end{promptbox}

\textbf{Step 2: 結果を見て追加指示}
\begin{promptbox}
ありがとうございます。以下の改善をお願いします：
1. 正規性の検定（Shapiro-Wilk検定）を先に行い、結果に応じて
   t検定またはMann-Whitney U検定を選択する
2. 効果量（Cohen's d）も計算する
3. 結果を辞書形式で返す
\end{promptbox}

\textbf{Step 3: さらに改善}
\begin{promptbox}
結果の辞書に以下も追加してください：
- 検定名（使用した検定の名前）
- 自由度（該当する場合）
- 95\%信頼区間
\end{promptbox}

このように、段階的に要件を追加することで、自分の研究に最適なコードに近づけていく。

%======================================================================
\section{実践テクニック}
\label{sec:6.4}

\subsection{段階的開発: 関数単位で生成・検証}

AIにコード全体を一度に生成させるのではなく、\textbf{関数単位で生成し、一つずつ動作確認する}方法が最も効率的である。

\begin{verbatim}
手順:
1. まず、プログラム全体の構成を相談する
2. 各関数を一つずつ生成させる
3. 生成された関数を実行して動作確認する
4. 問題があれば修正を依頼する
5. 全関数が揃ったら統合テストを行う
\end{verbatim}

\textbf{プロンプト例:}

\begin{promptbox}
以下の研究データ分析パイプラインを設計したいです。
まず、全体の構成（関数の一覧と役割）を提案してください。
コードはまだ書かなくて良いです。

処理の流れ:
1. 複数のCSVファイルを結合
2. データのクリーニング
3. 記述統計量の算出
4. 群間比較（統計検定）
5. 結果のグラフ作成
6. レポートの自動生成
\end{promptbox}

\subsection{テスト駆動: テストケースを先に書かせる}
\index{てすとくどうかいはつ@テスト駆動開発}

テスト駆動開発（TDD）の考え方をAIと組み合わせると、コードの品質向上が期待できる。

\textbf{手順:}

\begin{verbatim}
1. 関数の仕様を決める
2. AIにテストケースを先に生成させる
3. AIに関数本体を生成させる
4. テストを実行して合格を確認する
\end{verbatim}

\textbf{プロンプト例:}

\begin{promptbox}
以下の仕様の関数について、pytestのテストケースを先に書いてください。
関数本体はまだ書かないでください。

関数名: calculate\_effect\_size
入力: group1 (list[float]), group2 (list[float])
出力: dict（cohen\_d, interpretation を含む）
仕様:
- Cohen's d を計算する
- 解釈: |d| < 0.2 → "negligible", < 0.5 → "small",
  < 0.8 → "medium", >= 0.8 → "large"
- 空リストが渡された場合は ValueError を発生
\end{promptbox}

\subsection{リファクタリング依頼}
\index{りふぁくたりんぐ@リファクタリング}

既存のコードをAIに改善させることも有用な活用法である。

\textbf{プロンプト例:}

\begin{promptbox}
以下のPythonコードをリファクタリングしてください。
改善ポイント:
1. 可読性の向上（変数名の改善、コメント追加）
2. 関数への分割（1関数1責任の原則）
3. エラー処理の追加
4. 型ヒントの追加

{[}既存のコードをここに貼り付ける]
\end{promptbox}

\subsection{デバッグ支援: エラーメッセージの活用}
\index{でばっぐ@デバッグ}

エラーに遭遇したとき、AIは頼りになるデバッグ支援ツールとなる。

\textbf{プロンプト例:}

\begin{promptbox}
以下のPythonコードを実行したところ、エラーが発生しました。
原因と修正方法を教えてください。

\#\# コード
{[}コードを貼り付け]

\#\# エラーメッセージ
Traceback (most recent call last):
  File "analysis.py", line 42, in <module>
    result = df.groupby("condition").apply(analyze\_group)
ValueError: unequal length arrays

\#\# 追加情報
- Python 3.10, pandas 2.0, scipy 1.11
- データは条件ごとに被験者数が異なります
\end{promptbox}

\textbf{Tips:} エラーメッセージだけでなく、コード全体と実行環境の情報を含めることで、AIの回答精度が大幅に向上する。

%======================================================================
\section{テストコードの自動生成}
\label{sec:6.5}
\index{てすとこーど@テストコード}

\subsection{pytestによるテストケース生成}
\index{pytest@pytest}

テストコードの作成は、研究コードの信頼性を担保する上で極めて重要である。AIを活用すれば、包括的なテストケースを短時間で生成できる。

\textbf{プロンプト例:}

\begin{promptbox}
以下のPython関数に対するpytestのテストケースを生成してください。
正常系、異常系、境界値テストを含めてください。

{[}関数のコードを貼り付け]
\end{promptbox}

\textbf{対象関数の例:}

\begin{lstlisting}[language=Python]
def normalize_scores(
    scores: list[float],
    method: str = "min-max"
) -> list[float]:
    """スコアのリストを正規化する。"""
    if not scores:
        raise ValueError("スコアリストが空です")
    if method == "min-max":
        min_val, max_val = min(scores), max(scores)
        if min_val == max_val:
            return [0.0] * len(scores)
        return [(x - min_val) / (max_val - min_val)
                for x in scores]
    elif method == "z-score":
        mean_val = sum(scores) / len(scores)
        std_val = (sum((x - mean_val) ** 2
                   for x in scores) / len(scores)) ** 0.5
        if std_val == 0:
            return [0.0] * len(scores)
        return [(x - mean_val) / std_val for x in scores]
    else:
        raise ValueError(f"未対応の正規化方法: {method}")
\end{lstlisting}

\textbf{AIが生成するテストコードの例:}

\begin{lstlisting}[language=Python]
import pytest
from normalize import normalize_scores


class TestNormalizeScoresMinMax:
    """min-max正規化のテスト"""

    def test_basic_normalization(self):
        """基本的な正規化が正しく動作すること"""
        result = normalize_scores([0, 50, 100])
        assert result == [0.0, 0.5, 1.0]

    def test_negative_values(self):
        """負の値を含む場合"""
        result = normalize_scores([-10, 0, 10])
        assert result == [0.0, 0.5, 1.0]

    def test_single_element(self):
        """要素が1つの場合（min == max）"""
        result = normalize_scores([42.0])
        assert result == [0.0]

    def test_all_same_values(self):
        """全て同じ値の場合"""
        result = normalize_scores([5.0, 5.0, 5.0])
        assert result == [0.0, 0.0, 0.0]


class TestNormalizeScoresZScore:
    """z-score正規化のテスト"""

    def test_basic_z_normalization(self):
        """基本的なz-score正規化"""
        result = normalize_scores(
            [2, 4, 6], method="z-score"
        )
        assert pytest.approx(
            result[1], abs=1e-6
        ) == 0.0


class TestNormalizeScoresErrors:
    """異常系テスト"""

    def test_empty_list(self):
        """空リストでValueError"""
        with pytest.raises(
            ValueError, match="スコアリストが空です"
        ):
            normalize_scores([])

    def test_invalid_method(self):
        """不正なメソッド名でValueError"""
        with pytest.raises(
            ValueError, match="未対応の正規化方法"
        ):
            normalize_scores([1, 2, 3], method="invalid")


class TestNormalizeScoresBoundary:
    """境界値テスト"""

    def test_very_large_values(self):
        """非常に大きな値"""
        result = normalize_scores([0, 1e10])
        assert result == [0.0, 1.0]

    def test_two_elements(self):
        """要素が2つの最小ケース"""
        result = normalize_scores([10, 20])
        assert result == [0.0, 1.0]
\end{lstlisting}

\subsection{テストカバレッジの考え方}
\index{てすとかばれっじ@テストカバレッジ}

テストカバレッジとは、テストがコードのどの程度をカバーしているかを示す指標である。

\begin{lstlisting}[language=bash]
# カバレッジの測定
pip install pytest-cov
pytest --cov=my_module --cov-report=html tests/
\end{lstlisting}

研究コードでは、最低でも以下をカバーすることを目指す。

\begin{itemize}
\item \textbf{正常系}: 期待通りの入力に対する動作
\item \textbf{異常系}: 不正な入力に対するエラー処理
\item \textbf{境界値}: 最小値、最大値、空データ、要素1つなどの特殊なケース
\end{itemize}

\subsection{AI生成テストの検証ポイント}

AIが生成したテストコードは、以下の点を必ず確認する。

\begin{enumerate}
\item \textbf{テストが実際に実行できるか} --- インポートパスやファイル名が正しいか
\item \textbf{期待値が正しいか} --- AIが計算した期待値が本当に合っているか手計算で確認する
\item \textbf{テストの独立性} --- 各テストが他のテストに依存していないか
\item \textbf{意味のあるテストか} --- \texttt{assert True} のような無意味なテストが含まれていないか
\end{enumerate}

%======================================================================
\section{AIコードレビューとリファクタリング}
\label{sec:6.6}
\index{こーどれびゅー@コードレビュー}

\subsection{AIにコードレビューを依頼するプロンプト}

AIはコードレビューの補助ツールとしても活用できる。特に一人で研究を進める大学院生にとって、客観的なフィードバックを得る貴重な手段である。

\textbf{プロンプト例:}

\begin{promptbox}
以下のPythonコードをレビューしてください。
以下の観点でチェックし、改善提案をお願いします。

1. **可読性**: 変数名、関数名は適切か。コメントは十分か。
2. **効率性**: 不要なループや非効率な処理はないか。
3. **セキュリティ**: 入力検証は適切か。脆弱性はないか。
4. **エラー処理**: 例外処理は適切か。エッジケースへの対応は十分か。
5. **保守性**: コードの構造は適切か。将来の変更に対応しやすいか。

問題の深刻度を「重大」「中程度」「軽微」の3段階で分類してください。

{[}コードをここに貼り付け]
\end{promptbox}

\subsection{コードレビューチェックリスト}

AIにレビューを依頼する際の観点をチェックリストとしてまとめる。

\begin{itembox}[l]{コードレビューチェックリスト}
\textbf{可読性}
\begin{itemize}
\item 変数名・関数名が意図を表しているか
\item 適切なコメント・docstringがあるか
\item コードの構造が論理的か
\end{itemize}

\textbf{効率性}
\begin{itemize}
\item 不要なループや計算がないか
\item 適切なデータ構造を使っているか
\item メモリ使用量は適切か
\end{itemize}

\textbf{セキュリティ}
\begin{itemize}
\item ユーザー入力を適切に検証しているか
\item ファイルパスの検証があるか
\item SQLインジェクション等の脆弱性がないか
\end{itemize}

\textbf{エラー処理}
\begin{itemize}
\item 想定されるエラーを適切にキャッチしているか
\item エラーメッセージが分かりやすいか
\item リソースの解放が確実に行われるか
\end{itemize}

\textbf{テスト可能性}
\begin{itemize}
\item 関数が適切な粒度に分割されているか
\item 外部依存を注入可能か
\item 副作用が最小限か
\end{itemize}
\end{itembox}

\subsection{リファクタリングの種類と適用場面}

\begin{table}[htbp]
\centering
\caption{リファクタリングの種類と適用場面}
\label{tab:refactoring}
\footnotesize
\begin{tabular}{lll}
\toprule
種類 & 適用場面 & プロンプト例 \\
\midrule
関数の抽出 & 長い関数を分割 & 「この関数を責務ごとに \\
 & & 小さな関数に分割して」 \\
変数名の改善 & 意味不明な変数名 & 「変数名をより意味が \\
 & & 通るように改善して」 \\
重複の排除 & 類似コードの繰り返し & 「重複しているコードを \\
 & & 共通関数にまとめて」 \\
条件分岐の簡略化 & 複雑なif-else & 「この条件分岐を \\
 & & 簡潔に書き直して」 \\
データ構造の変更 & 不適切な構造 & 「リストの代わりに辞書を \\
 & & 使うように変更して」 \\
\bottomrule
\end{tabular}
\end{table}

\subsection{ドキュメント自動生成}

AIを使って、コードのドキュメントを自動生成できる。

\textbf{docstring生成プロンプト:}

\begin{promptbox}
以下のPython関数にNumPyスタイルのdocstringを追加してください。
Parameters、Returns、Raises、Examples セクションを含めてください。

{[}関数のコードを貼り付け]
\end{promptbox}

\textbf{README生成プロンプト:}

\begin{promptbox}
以下のPythonプロジェクトのREADME.mdを作成してください。

プロジェクト名: eeg-analysis-toolkit
概要: EEGデータの前処理・分析パイプライン
主要ファイル:
- preprocess.py: データ前処理
- analyze.py: 統計分析
- visualize.py: 結果の可視化
- requirements.txt: 依存パッケージ

以下のセクションを含めてください:
1. プロジェクト概要
2. インストール方法
3. 使い方（コード例付き）
4. ファイル構成
5. ライセンス
\end{promptbox}

%======================================================================
\section{AI生成コードの品質と注意点}
\label{sec:6.7}
\index{せきゅりてぃ@セキュリティ}

\subsection{セキュリティリスク}

AI生成コードには、以下のようなセキュリティリスクが含まれる場合がある。

\textbf{1. SQLインジェクション}

\begin{lstlisting}[language=Python]
# 危険: AI が生成しがちなコード
def get_user(name):
    query = f"SELECT * FROM users WHERE name = '{name}'"
    cursor.execute(query)

# 安全: パラメータ化クエリを使用
def get_user(name):
    query = "SELECT * FROM users WHERE name = ?"
    cursor.execute(query, (name,))
\end{lstlisting}

\textbf{2. 不適切な入力検証}

\begin{lstlisting}[language=Python]
# 危険: ファイルパスを検証していない
def read_file(path):
    with open(path) as f:
        return f.read()

# 安全: パスの検証を含む
def read_file(path, allowed_dir="/data/experiments"):
    real_path = os.path.realpath(path)
    if not real_path.startswith(
        os.path.realpath(allowed_dir)
    ):
        raise ValueError("許可されていないパスです")
    with open(real_path) as f:
        return f.read()
\end{lstlisting}

\textbf{3. ハードコードされた秘密情報}

\begin{lstlisting}[language=Python]
# 危険: APIキーがコードに含まれている
api_key = "sk-1234567890abcdef"

# 安全: 環境変数から読み込む
import os
api_key = os.environ.get("API_KEY")
if not api_key:
    raise RuntimeError(
        "API_KEY 環境変数が設定されていません"
    )
\end{lstlisting}

\subsection{ライセンス問題}
\index{らいせんす@ライセンス}

AIが生成したコードには、学習データに含まれるオープンソースコードの断片が含まれる可能性がある。

\begin{itemize}
\item \textbf{研究コードの場合}: 論文の再現性のために公開する場合、ライセンスを明記する
\item \textbf{推奨}: MITライセンスやApache 2.0ライセンスなど、許容的なライセンスを選択する
\item \textbf{注意}: AI生成コードの著作権は法的にグレーゾーンであることを認識しておく
\end{itemize}

\subsection{「信頼するが検証する」の原則}

AI生成コードに対する基本姿勢は\textbf{「Trust, but verify（信頼するが検証する）」}である。

\begin{itembox}[l]{検証チェックリスト}
\begin{itemize}
\item コードが意図通りに動作するか（手動テスト）
\item テストケースが全て通過するか（自動テスト）
\item エッジケースでも正しく動作するか
\item エラー処理が適切か
\item パフォーマンスが許容範囲内か
\item セキュリティ上の問題がないか
\item コードの論理が自分で理解できるか $\leftarrow$ \textbf{最も重要}
\end{itemize}
\end{itembox}

最後の項目が最も重要である。\textbf{自分で理解できないコードは使わない}。AIに説明を求め、納得してから使用すること。

\subsection{研究コードの再現性確保}
\index{さいげんせい@再現性}

研究で使用するコードは、再現性が欠かせない。以下の対策を講じる。

\begin{enumerate}
\item \textbf{依存関係の固定}: \texttt{requirements.txt}や\texttt{pyproject.toml}でバージョンを固定する
\item \textbf{乱数シードの設定}: 結果の再現性のために、乱数シードを明示的に設定する
\item \textbf{環境情報の記録}: Pythonバージョン、OS、主要ライブラリのバージョンを記録する
\item \textbf{バージョン管理}: Gitでコードの履歴を管理する（第8章で詳述）
\end{enumerate}

\begin{lstlisting}[language=Python]
# 再現性確保のテンプレート
import random
import numpy as np

RANDOM_SEED = 42
random.seed(RANDOM_SEED)
np.random.seed(RANDOM_SEED)

# 実行環境の記録
import sys
import platform
print(f"Python: {sys.version}")
print(f"NumPy: {np.__version__}")
print(f"OS: {platform.platform()}")
\end{lstlisting}

%======================================================================
\section*{演習問題}

\begin{enumerate}

\item \textbf{演習6-1: 基本的なコード生成}

以下の仕様で、AIを使ってPython関数を生成しなさい。生成されたコードを実行し、正しく動作することを確認しなさい。

\textbf{仕様:}
\begin{itemize}
\item 関数名: \texttt{summarize\_data}
\item 入力: pandasのDataFrame
\item 出力: 各数値列について、平均、中央値、標準偏差、最小値、最大値を含む辞書
\item エラー処理: DataFrameが空の場合はValueErrorを発生
\end{itemize}

使用したプロンプトと、生成されたコードを提出すること。

\item \textbf{演習6-2: テスト駆動開発}

演習6-1で作成した\texttt{summarize\_data}関数に対して、AIを使ってpytestのテストケースを生成しなさい。以下のケースを最低限含めること。

\begin{itemize}
\item 正常なDataFrameに対するテスト
\item 空のDataFrameに対するテスト
\item 数値列のないDataFrameに対するテスト
\item 欠損値を含むDataFrameに対するテスト
\end{itemize}

\item \textbf{演習6-3: コードレビューとリファクタリング}

以下のコードをAIにレビューさせ、改善版を作成しなさい。どのような問題が指摘され、どう改善されたかをレポートにまとめること。

\begin{lstlisting}[language=Python]
def proc(d):
    r = []
    for i in range(len(d)):
        if d[i] != None and d[i] > 0:
            x = d[i] * 2.5 + 10
            r.append(x)
    s = 0
    for j in r:
        s = s + j
    a = s / len(r)
    return a
\end{lstlisting}

\item \textbf{演習6-4: デバッグ支援}

以下のコードにはバグが含まれている。AIを使ってバグを特定し、修正しなさい。修正の過程で使用したプロンプトも提出すること。

\begin{lstlisting}[language=Python]
import pandas as pd

def calculate_growth_rate(df, value_col, time_col):
    """時系列データの成長率を計算する"""
    df = df.sort_values(time_col)
    growth_rates = []
    for i in range(len(df)):
        current = df.iloc[i][value_col]
        previous = df.iloc[i-1][value_col]
        rate = (current - previous) / previous * 100
        growth_rates.append(rate)
    df["growth_rate"] = growth_rates
    return df

# テストデータ
data = pd.DataFrame({
    "year": [2020, 2021, 2022, 2023],
    "revenue": [100, 120, 150, 180]
})

result = calculate_growth_rate(data, "revenue", "year")
print(result)
\end{lstlisting}

\item \textbf{演習6-5: 研究データ分析スクリプトの作成}

自分の研究テーマ（または以下のサンプルテーマ）に関連するデータ分析スクリプトをAIと対話的に作成しなさい。

\textbf{サンプルテーマ}: 「オンライン授業と対面授業の学習効果比較」

以下の機能を含むスクリプトを作成すること。
\begin{enumerate}
\item サンプルデータの生成（被験者50名、2群、テストスコアあり）
\item 記述統計量の算出
\item 適切な統計検定の実施
\item 結果のグラフ作成
\end{enumerate}

AIとの対話履歴（プロンプトとレスポンスの要約）、最終的なコード、実行結果を提出すること。

\end{enumerate}

\vspace{10mm}
\begin{itembox}[l]{本章のまとめ}
AIは頼りになるコーディングパートナーであるが、最終的な品質責任は使用者にある。「AIに書かせて、人間が検証する」という姿勢を常に持ち、特に研究コードでは再現性とセキュリティに十分注意すること。プロンプトの書き方一つでコードの品質が大きく変わるため、明確で具体的な仕様記述のスキルを磨くことが重要である。
\end{itembox}
